\section{\S 4.2 原电池}\label{sect:原电池}

    \subsection{Zn-Cu 原电池}\label{subsec:ZnCu原电池}

        \begin{center}
        % 图：Zn-Cu 原电池（稀硫酸，单液；Zn－，Cu＋，外接电流表）
\begin{tikzpicture}[x=1cm, y=1cm, every node/.style={font=\small\ssans}]
    % 烧杯
    \draw[rounded corners=8pt] (-2.1,0) -- (-2.1,-3.3) -- (2.1,-3.3) -- (2.1,0);
    % 稀硫酸
    \draw[rounded corners=8pt, fill=blue!8] (-2.1,-3.3) -- (-2.1,-1.0) -- (2.1,-1.0) -- (2.1,-3.3) -- cycle;
    % 电极
    \draw[fill=gray!55, draw=black!70] (-1.7,-1.05) rectangle (-1.45,-2.95);
    \draw[fill=black!60, draw=black!80] (1.45,-1.05) rectangle (1.7,-2.95);
    % 导线与电流表
    \draw (-1.575,-1.05) -- (-1.575,0.9);
    \draw (1.575,-1.05) -- (1.575,0.9);
    \draw (-1.575,0.9) -- (-0.5,0.9) (0.5,0.9) -- (1.575,0.9);
    \draw[fill=white] (-0.5,0.9) circle (0.45);
    \node at (0,0.9) {A};
    % 电子流向（外电路）
    \draw[->,blue!60!black,very thick] (-1.35,1.35) -- (1.2,1.35);
    \node[blue!60!black] at (0,1.62) {\cemh{e-}};
    % 标注
    \node[left=1pt] at (-1.575,-1.6) {Zn（\textcolor{blue}{$-$}）};
    \node[right=1pt] at (1.575,-1.6) {Cu（\textcolor{blue}{$+$}）};
    \node at (0,-3.0) {稀硫酸};
    \node[font=\small] at (0,-3.7) {\cemh{Zn - 2e- = Zn^{2+}}\quad \cemh{2H+ + 2e- = H2 ^}};
    % 电极反应离子箭头（溶液中）
    \draw[->,gray!70,thick] (-1.575,-2.6) -- (-1.15,-1.35) node[midway,left,font=\scriptsize] {Zn\textsuperscript{2+}};
    \draw[->,gray!70,thick] (1.575,-2.7) -- (1.2,-1.4) node[midway,right,font=\scriptsize] {H\textsuperscript{+}};
\end{tikzpicture}

        \end{center}

        \cemh{Cu}：\textbf{正极}

        \cemh{Zn}：\textbf{负极}

        氧化：\cemh{Zn - 2e- \xleq{} Zn^{2+}}

        还原：\cemh{2H+ + 2e- \xleq{} H2 ^}

        \textcolor{red}{正正负负}\(\left\{\begin{matrix}
            \text{正离子} \to \text{正极}\\
            \text{负离子} \to \text{负极}
        \end{matrix}\right.\)

        电子：负极（\cemh{Zn}）\(\to\) 正极（\cemh{Cu}）（导线）

        盐桥：离子导体，平衡电荷

    \subsection{构成原电池的条件}\label{subsec:构成原电池条件}

        \begin{enumerate}
            \item 两个活泼性不同的能导电的电极：\Circled{1} 活泼的作负极、\Circled{2} 不活泼的作正极
            \item 电极间形成闭合回路（电子导体）
            \item 具有电解质溶液或熔融电解质（离子导体）
            \item 自发的氧化还原反应
        \end{enumerate}

        \cemh{Zn/Cu/ZnSO4} 溶液 \cmark

        正极：\cemh{O2 + 2H2O + 4e- \xleq{} 4OH-}

        \cemh{Fe/C/NaCl} 溶液 \cmark

        \textcolor{blue}{水中溶解的 \cemh{O2} 可以自发氧化大多数金属。}

    \subsection{单液原电池与双液原电池}\label{subsec:单液双液原电池}

        \noindent
        \begin{minipage}[t]{0.48\textwidth}
            \centering
            单液原电池\\[2pt]
            \resizebox{0.98\linewidth}{!}{% 图：单液原电池（Zn 与 Cu²⁺ 直接接触，电流不稳定）
\begin{tikzpicture}[x=1cm, y=1cm, every node/.style={font=\small\ssans}]
    % 烧杯与 CuSO4 溶液
    \draw[rounded corners=8pt] (-2.1,0) -- (-2.1,-3.3) -- (2.1,-3.3) -- (2.1,0);
    \draw[rounded corners=8pt, fill=blue!18] (-2.1,-3.3) -- (-2.1,-1.0) -- (2.1,-1.0) -- (2.1,-3.3) -- cycle;
    % 电极（单液：Zn、Cu 同液，Zn 表面直接置换）
    \draw[fill=gray!55, draw=black!70] (-1.7,-1.05) rectangle (-1.45,-2.95);
    \draw[fill=black!60, draw=black!80] (1.45,-1.05) rectangle (1.7,-2.95);
    \draw[fill=black!60, draw=black!80] (-1.1,-2.95) rectangle (-0.85,-2.5);
    % 导线与电流表
    \draw (-1.575,-1.05) -- (-1.575,0.9);
    \draw (1.575,-1.05) -- (1.575,0.9);
    \draw (-1.575,0.9) -- (-0.5,0.9) (0.5,0.9) -- (1.575,0.9);
    \draw[fill=white] (-0.5,0.9) circle (0.45);
    \node at (0,0.9) {A};
    \draw[->,blue!60!black,very thick] (-1.35,1.35) -- (1.2,1.35);
    \node[blue!60!black] at (0,1.62) {\cemh{e-}};
    % 标注
    \node[left=1pt] at (-1.575,-1.6) {Zn（\textcolor{blue}{$-$}）};
    \node[right=1pt] at (1.575,-1.6) {Cu（\textcolor{blue}{$+$}）};
    \node[font=\scriptsize, text=black!70] at (0,-1.35) {\cemh{Zn + Cu^{2+}} 直接置换};
    \node[font=\scriptsize, gray!60!black] at (-1.0,-2.75) {置换出的 Cu};
    \node at (0,-3.0) {\cemh{CuSO4} 溶液};
    \node[font=\small\ssans, text=red] at (0,-3.7) {Zn 与 \cemh{Cu^{2+}} 直接接触，电流不稳定};
\end{tikzpicture}
}
        \end{minipage}
        \hfill
        \begin{minipage}[t]{0.48\textwidth}
            \centering
            双液原电池\\[2pt]
            \resizebox{0.98\linewidth}{!}{% 图：双液原电池（盐桥连接 ZnSO4 / CuSO4 两半电池）
\begin{tikzpicture}[x=1cm, y=1cm, every node/.style={font=\small\ssans}]
    % 两个烧杯
    \draw[rounded corners=6pt] (-3.5,0) -- (-3.5,-3.2) -- (-1.2,-3.2) -- (-1.2,0);
    \draw[rounded corners=6pt] (1.2,0) -- (1.2,-3.2) -- (3.5,-3.2) -- (3.5,0);
    % 溶液
    \draw[rounded corners=6pt, fill=blue!8] (-3.5,-3.2) -- (-3.5,-1.0) -- (-1.2,-1.0) -- (-1.2,-3.2) -- cycle;
    \draw[rounded corners=6pt, fill=blue!20] (1.2,-3.2) -- (1.2,-1.0) -- (3.5,-1.0) -- (3.5,-3.2) -- cycle;
    % 电极
    \draw[fill=gray!55, draw=black!70] (-3.0,-1.05) rectangle (-2.75,-2.9);
    \draw[fill=black!60, draw=black!80] (2.75,-1.05) rectangle (3.0,-2.9);
    % 盐桥（倒 U 管，含 KNO3 凝胶）
    \draw[thick, fill=white] (-1.35,-0.9) .. controls (-1.35,0.55) and (1.35,0.55) .. (1.35,-0.9);
    \draw[thick, fill=blue!5] (-1.35,-0.9) .. controls (-1.35,0.4) and (1.35,0.4) .. (1.35,-0.9);
    \node[font=\scriptsize] at (0,0.32) {盐桥（\cemh{KNO3}）};
    % 导线与电流表
    \draw (-2.875,-1.05) -- (-2.875,1.5);
    \draw (2.875,-1.05) -- (2.875,1.5);
    \draw (-2.875,1.5) -- (-0.5,1.5) (0.5,1.5) -- (2.875,1.5);
    \draw[fill=white] (-0.5,1.5) circle (0.45);
    \node at (0,1.5) {A};
    \draw[->,blue!60!black,very thick] (-2.6,1.95) -- (2.55,1.95);
    \node[blue!60!black] at (0,2.25) {\cemh{e-}};
    % 离子迁移方向
    \draw[->,gray!60!black,thick] (-3.1,-2.2) -- (-3.1,-1.35);
    \draw[->,gray!60!black,thick] (3.1,-1.35) -- (3.1,-2.2);
    \node[font=\scriptsize, gray!60!black, left=0pt] at (-3.1,-1.9) {\cemh{Zn^{2+}}\ 进\ 入\ 溶\ 液};
    \node[font=\scriptsize, gray!60!black, right=0pt] at (3.1,-1.9) {\cemh{Cu^{2+}}\ 放\ 电\ 析\ 出};
    % 标注
    \node at (-2.375,-3.0) {\cemh{ZnSO4} 溶液};
    \node at (2.375,-3.0) {\cemh{CuSO4} 溶液};
    \node[left=1pt] at (-2.875,-1.6) {Zn（\textcolor{blue}{$-$}）};
    \node[right=1pt] at (2.875,-1.6) {Cu（\textcolor{blue}{$+$}）};
\end{tikzpicture}
}
        \end{minipage}

        基于 \cemh{Zn + Cu^{2+} \xleq{} Cu + Zn^{2+}} 设计双液原电池：

        氧化：\cemh{Zn - 2e- \xleq{} Zn^{2+}}

        还原：\cemh{Cu^{2+} + 2e- \xleq{} Cu}

    \subsection{离子交换膜}\label{subsec:离子交换膜}

        \begin{enumerate}
            \item 正离子交换膜：只允许正离子
            \item 负离子交换膜：只允许负离子
            \item 质子交换膜：只允许 \cemh{H+}
        \end{enumerate}

        \textcolor{blue}{\cemh{H+} 在正极得 \cemh{e-} 无阻更快}

        \textcolor{blue}{纯 \cemh{Zn}：理论无气泡}

        \textcolor{blue}{粗 \cemh{Zn}：速率快}

        \textcolor{red}{电极 = 电极材料 + 与电极接触的物质}

    \subsection{浓差电池}\label{subsec:浓差电池}

        浓差电池：利用浓淡差。

        金属失电子难易与溶液浓度相关。

        % TODO: 来源待补充 ——【浓差电池】内容已删除（原 content 来源未知），如需可依据课本重写
